\section{Memory Fences}

\textbf{What are memory fences?}

In JavaScript/TypeScript, when you write \texttt{array[i] = 42;}, it happens immediately and in order. But in GPUs and CPUs, the hardware can reorder memory operations for performance. A \textbf{memory fence} is like saying ``wait, make sure all pending writes are actually visible before continuing.''

\textbf{Barriers vs. Fences:}

Think of them like this:
\begin{itemize}[noitemsep]
  \item \textbf{Barrier} (\texttt{syncThreads()}): ``Everyone stop and wait here. Don't proceed until all threads arrive.'' (Synchronization + memory ordering)
  \item \textbf{Fence} (\texttt{threadfence()}): ``Make my writes visible, but don't wait for other threads.'' (Memory ordering only)
\end{itemize}

\textbf{When to use fences:}

Use fences for lock-free algorithms and atomic operations where you need memory ordering guarantees but don't want the overhead of waiting for all threads.

\textbf{Example: \acr{GPU} memory fences}

\begin{lstlisting}[style=ts]
import { FenceInst, CallInst } from "@euriklis/llvm-ir";

// Standard LLVM fence (works on all architectures)
const fence = new FenceInst({
  ordering: "seq_cst"  // Sequential consistency (strongest)
});
// Generates: fence seq_cst

// CUDA threadfence (block-level) - ensures writes visible to block
const fenceBlock = new CallInst({
  functionName: "llvm.nvvm.membar.cta",
  returnType: voidType,
  args: []
});
// Generates: call void @llvm.nvvm.membar.cta()

// CUDA threadfence (global-level) - ensures writes visible globally
const fenceGlobal = new CallInst({
  functionName: "llvm.nvvm.membar.gl",
  returnType: voidType,
  args: []
});
// Generates: call void @llvm.nvvm.membar.gl()
\end{lstlisting}

\textbf{Memory ordering levels:}

\texttt{FenceInst} supports different ordering strengths (from weakest to strongest):
\begin{itemize}[noitemsep]
  \item \texttt{"acquire"}: Ensures reads after fence see all writes before
  \item \texttt{"release"}: Ensures writes before fence are visible after
  \item \texttt{"acq\_rel"}: Both acquire and release
  \item \texttt{"seq\_cst"}: Sequentially consistent (default, safest)
\end{itemize}

For most use cases, use \texttt{"seq\_cst"} (sequential consistency). Only use weaker orderings if you understand low-level memory models and need the performance.

